%%
%% spring-lecture10.tex
%% 
%% Made by Alex Nelson <pqnelson@gmail.com>
%% Login   <alex@lisp>
%% 
%% Started on  2026-04-23T12:27:54-0700
%% Last update 2026-04-23T12:27:54-0700
%% 

\lecture[Geodesics]{}

\begin{node}
In Euclidean space, to connect two points $A$ and $B$ in a path of
minimal length, we just take a straight line. But on a manifold,
``straight line'' is no longer a sensible notion. However,
\emph{curves} are a sensible notion, as is ``minimal length''.

If there always exists
at least one such curve for any two points on a manifold, then we call
it ``complete''.

So on Euclidean space $(\RR^{n},g_{0})$ we always have a curve
\begin{equation}
\gamma\colon[0,1]\to\RR^{n}
\end{equation}
such that $\gamma(0)=A$ and $\gamma(1)=B$. Moreover, $\gamma$ is a
path of minimum length if and only if $\gamma$ is a straight line
(``if and only if'' $\gamma'(t)=\gamma(1)-\gamma(0)$ for all $t\in[0,1]$).
Saying a curve ``has constant velocity vector'' makes sense now:
$\gamma'(t)$ is a vector field on the curve $\gamma$.
\end{node}

\begin{definition}
Let $(M,g)$ be a Riemannian manfiold. Let $I\subset\RR$ be an
interval. Let $\gamma\colon I\to M$ be a smooth curve. We say $\gamma$
is a \define{Geodesic} if its acceleration vanishes identically:
\begin{equation}
\frac{D\gamma'}{\D t}=0.
\end{equation}
(This implicitly depends on the connection by relying on $D/\D t$, we
have decided to set it to the Levi-Civita connection by convention.)
\end{definition}

\begin{remark}[Form of geodesic equations in a local chart]
Let $(M,g)$ be a Riemannian manifold. Let
$(U,\varphi=(\xi^{1},\dots,\xi^{n}))\in\smoothstr{M}$ be a local
chart. Consider a curve
\begin{equation}
\gamma\colon I\to M
\end{equation}
such that $\gamma(I)\subset U$. Then
\begin{equation}
\frac{D\gamma'}{\D t}=\left.\sum^{n}_{k=1}((\gamma^{k})''(t)+\sum^{n}_{i,j=1}\Gamma^{k}_{ij}(\gamma^{i})'(t)(\gamma^{j})'(t))\frac{\partial}{\partial\xi^{k}}\right|_{\gamma(t)}
\end{equation}
where $\gamma^{i}=(\xi^{i}\circ\gamma)\colon I\to\RR$ for every $i=1,\dots,n$.
Here
\begin{equation}
\Gamma^{k}_{ij}=\Gamma^{k}_{ij}(\gamma(t)).
\end{equation}
Therefore $\gamma$ isa geodesic if and only if
\begin{equation}
0=\left.\sum^{n}_{k=1}((\gamma^{k})''(t)+\sum^{n}_{i,j=1}\Gamma^{k}_{ij}(\gamma^{i})'(t)(\gamma^{j})'(t))\frac{\partial}{\partial\xi^{k}}\right|_{\gamma(t)}
\end{equation}
for every $k=1,\dots,n$ and every $t\in I$.
\end{remark}

\begin{theorem}
Let $(M,g)$ be a Riemannian manifold. Then for every $x_{0}\in M$,
there exists an open neighborhood $U\subset M$ of $x_{0}\in U$ and
there exists an $\varepsilon>0$ such that for every $(x,v)\in TU$ with
$|v|_{g}<\varepsilon$, there exists a unique geodesic
\begin{equation}
\alpha_{x,v}\colon(-1,1)\to M
\end{equation}
such that
\begin{enumerate}
\item $\alpha_{x,v}(0)=x$ and
\item $\alpha_{x,v}'(0)=v$.
\end{enumerate}
Moreover, the map $\psi\colon TU\times(-1,1)\to M$ given by
\begin{equation}
\psi\bigl((x,v),t\bigr)=\alpha_{x,v}(t)
\end{equation}
is smooth.
\end{theorem}

\begin{remark}
If you find yourself ina situation when you found a geodesic on a
smaller interval, say $\alpha_{x,v}\colon(-1/2,1/2)\to M$, then you
just picked the wrong initial velocity $v$. To obtain a desirable
$\widetilde{\alpha}_{x,v}\colon(-1,1)\to M$ we just define it by
$t\mapsto\alpha_{x,v}(t/2)$ to get the desired result.
\end{remark}

\begin{remark}
If you have two geodesics that are both defined at time $t=0$, and
both geodesics have the same initial conditions, then they agree
everywhere they overlap.
\end{remark}

\begin{definition}[Maximal geodesic]
A geodesic $\gamma\colon I\to M$ with $0\in I$ is called
\define{Maximal} if for every
$\widetilde{\gamma}\colon\widetilde{I}\to M$ with $0\in\widetilde{I}$
such that they have the same initial conditions
$\widetilde{\gamma}(0)=\gamma(0)$ and $\widetilde{\gamma}'(0)=\gamma'(0)$
we have $\widetilde{I}\subset I$.
\end{definition}

\begin{remark}
Given $x\in M$ and $v\in T_{x}M$, there always exists a maximal
geodesic
\begin{equation}
\gamma_{x,v}\colon I\to M
\end{equation}
with $0\in I$ and initial conditions $\gamma_{x,v}(0)=x$ and
$\gamma'_{x,v}(0)=v$. (This interval $I$ could be small if $v$ is
chosen to be ``too big. Just pick $v$ with a smaller magnitude $v/N$.)
\end{remark}

\begin{remark}
Maximality is proven by Zorn's lemma.
\end{remark}

\begin{definition}[Exponential map]
Let $(M,g)$ be a Riemannian manifold. Cosnider the open set
$\Omega=\{(x,v)\in TM\mid \gamma_{x,v}\mbox{ exists at time }t=1\}$.
The \define{Exponential Map} $\exp\colon\Omega\to M$ is given by
\begin{equation}
\exp(x,v):=\gamma_{x,v}(1),
\end{equation}
for any $(x,v)\in\Omega$.

Note: this is a smooth map.
\end{definition}

\begin{notation}
We will switch notations, writing $\exp(x,v)$ as $\exp_{x}(v)$ when convenient.
\end{notation}

\begin{xca}
Show the differential of $\exp(x,0)$ is $\D(\exp_{x})_{0}=\id$.
\end{xca}

\begin{remark}
This exponential map originates from Lie groups. For example, take
$\SO(n)$. This is naturally a smooth submanifold of
$\RR^{n^{2}}$. We see
\begin{equation}
\begin{split}
T_{\id}\SO(n) &= \operatorname{Lie}(\SO(n))\\
&=\{A\in M_{n}(\RR)\mid A^{\mathsf{T}}=-A\}.
\end{split}
\end{equation}
Then $\exp_{\id}(B)$ is just the matrix exponential of $B\in T_{\id}\SO(n)$.
\end{remark}